\begin{titlepage}
\thispagestyle{empty}
\sffamily
\setlength{\parindent}{0pt}
\vspace*{0.03\textheight}

{\large\bfseries\color{DefinitionInk} A Guide to the Formal Proof in Lean\par}
\vspace{0.55em}
{\color{DefinitionInk}\rule{27mm}{1pt}\par}
\vspace{1.35em}

{\fontsize{29pt}{36pt}\linespread{1.06}\selectfont\bfseries\color{OxfordNavy}
 From Theta Lifts to the Proof\par
 of the Generalized Riemann\par
 Hypothesis\par}
\vspace{1.65em}

{\normalsize\rmfamily\color{SlateInk}
 Based on the readable Lean proof in
 \href{https://github.com/lsernau/grh-lean}{\texttt{lsernau/grh-lean}}
 and its accompanying blueprint.\par}
\vspace{0.35em}
{\small\sffamily\color{SlateInk!75}
 Upstream build environment: Lean 4.35.0\quad\textperiodcentered\quad
 Mathlib \texttt{56a00585}\par}

\vfill
{\graphicspath{{../figures/}}%
 \noindent% Shared by the Chinese and English covers.
% Tall borderless monochrome-gray editorial monograph plate with native TikZ critical-strip plot.
\providecommand{\CoverLink}[2]{\hyperref[#1]{#2}}
\ifcsname CoverSansFont\endcsname\else
  \newfontfamily\CoverSansFont{texgyreheros}[
    Extension=.otf, UprightFont=*-regular, ItalicFont=*-italic,
    BoldFont=*-bold, BoldItalicFont=*-bolditalic]
\fi
\ifcsname CoverSerifFont\endcsname\else
  \newfontfamily\CoverSerifFont{texgyrepagella}[
    Extension=.otf, UprightFont=*-regular, ItalicFont=*-italic,
    BoldFont=*-bold, BoldItalicFont=*-bolditalic]
\fi

% Soft monochrome architectural gray palette (no harsh black rules; DeepBurgundy only for zeros/RH)
\definecolor{CoverBg}{RGB}{244,245,247}
\definecolor{CoverStripLeft}{RGB}{231,234,238}
\definecolor{CoverStripRight}{RGB}{238,240,243}
\definecolor{CoverWaveFill}{RGB}{118,125,134}
\definecolor{CoverWaveLine}{RGB}{86,93,102}
\definecolor{CoverInk}{RGB}{50,55,62}
\definecolor{CoverMuted}{RGB}{98,105,114}
\definecolor{CoverLightInk}{RGB}{138,145,154}
\definecolor{CoverTierLight}{RGB}{156,163,172}
\definecolor{CoverTierDark}{RGB}{94,101,110}
\definecolor{CoverRule}{RGB}{204,209,216}
\definecolor{CoverHairline}{RGB}{218,222,228}
\definecolor{CoverGrid}{RGB}{217,221,227}
\definecolor{CoverAxis}{RGB}{108,115,124}

% Load mpmath-generated critical-strip zeros and Hardy Z(t) curve coordinates
\IfFileExists{figures/critical-strip-data.tex}{%
  % Auto-generated by plot_critical_strip.py using 50-digit mpmath arithmetic.
% Coordinates are in critical-strip units: x = Re(s), y = Im(s) on [0, 40].
\providecommand{\CoverCriticalZeros}{%
  14.134725, 21.022040, 25.010858, 30.424876, 32.935062, 37.586178%
}
\providecommand{\CoverFillHardyZArea}[1][]{%
  \fill[#1] (0.5,0) -- plot coordinates {
    (0.33206,0.00000) (0.33580,0.11111) (0.34557,0.22222) (0.35830,0.33333)
    (0.37134,0.44444) (0.38326,0.55556) (0.39355,0.66667) (0.40220,0.77778)
    (0.40939,0.88889) (0.41532,1.00000) (0.42022,1.11111) (0.42426,1.22222)
    (0.42759,1.33333) (0.43033,1.44444) (0.43258,1.55556) (0.43440,1.66667)
    (0.43587,1.77778) (0.43704,1.88889) (0.43794,2.00000) (0.43861,2.11111)
    (0.43908,2.22222) (0.43936,2.33333) (0.43948,2.44444) (0.43944,2.55556)
    (0.43928,2.66667) (0.43899,2.77778) (0.43858,2.88889) (0.43807,3.00000)
    (0.43745,3.11111) (0.43674,3.22222) (0.43594,3.33333) (0.43505,3.44444)
    (0.43408,3.55556) (0.43304,3.66667) (0.43191,3.77778) (0.43071,3.88889)
    (0.42944,4.00000) (0.42810,4.11111) (0.42669,4.22222) (0.42521,4.33333)
    (0.42367,4.44444) (0.42206,4.55556) (0.42039,4.66667) (0.41866,4.77778)
    (0.41688,4.88889) (0.41503,5.00000) (0.41313,5.11111) (0.41117,5.22222)
    (0.40916,5.33333) (0.40710,5.44444) (0.40499,5.55556) (0.40282,5.66667)
    (0.40062,5.77778) (0.39837,5.88889) (0.39607,6.00000) (0.39374,6.11111)
    (0.39137,6.22222) (0.38897,6.33333) (0.38653,6.44444) (0.38407,6.55556)
    (0.38158,6.66667) (0.37907,6.77778) (0.37655,6.88889) (0.37401,7.00000)
    (0.37146,7.11111) (0.36891,7.22222) (0.36635,7.33333) (0.36380,7.44444)
    (0.36127,7.55556) (0.35874,7.66667) (0.35624,7.77778) (0.35377,7.88889)
    (0.35133,8.00000) (0.34893,8.11111) (0.34658,8.22222) (0.34428,8.33333)
    (0.34205,8.44444) (0.33988,8.55556) (0.33779,8.66667) (0.33579,8.77778)
    (0.33388,8.88889) (0.33207,9.00000) (0.33037,9.11111) (0.32879,9.22222)
    (0.32734,9.33333) (0.32602,9.44444) (0.32485,9.55556) (0.32384,9.66667)
    (0.32300,9.77778) (0.32233,9.88889) (0.32184,10.00000) (0.32155,10.11111)
    (0.32147,10.22222) (0.32159,10.33333) (0.32195,10.44444) (0.32253,10.55556)
    (0.32336,10.66667) (0.32443,10.77778) (0.32577,10.88889) (0.32738,11.00000)
    (0.32926,11.11111) (0.33142,11.22222) (0.33388,11.33333) (0.33664,11.44444)
    (0.33970,11.55556) (0.34308,11.66667) (0.34677,11.77778) (0.35078,11.88889)
    (0.35511,12.00000) (0.35978,12.11111) (0.36477,12.22222) (0.37010,12.33333)
    (0.37576,12.44444) (0.38176,12.55556) (0.38808,12.66667) (0.39474,12.77778)
    (0.40172,12.88889) (0.40902,13.00000) (0.41664,13.11111) (0.42456,13.22222)
    (0.43279,13.33333) (0.44131,13.44444) (0.45011,13.55556) (0.45918,13.66667)
    (0.46850,13.77778) (0.47806,13.88889) (0.48785,14.00000) (0.49785,14.11111)
    (0.50000,14.13473) (0.50804,14.22222) (0.51839,14.33333) (0.52890,14.44444)
    (0.53952,14.55556) (0.55025,14.66667) (0.56106,14.77778) (0.57192,14.88889)
    (0.58279,15.00000) (0.59367,15.11111) (0.60450,15.22222) (0.61527,15.33333)
    (0.62595,15.44444) (0.63649,15.55556) (0.64687,15.66667) (0.65705,15.77778)
    (0.66701,15.88889) (0.67670,16.00000) (0.68610,16.11111) (0.69516,16.22222)
    (0.70385,16.33333) (0.71214,16.44444) (0.72000,16.55556) (0.72738,16.66667)
    (0.73427,16.77778) (0.74062,16.88889) (0.74641,17.00000) (0.75161,17.11111)
    (0.75618,17.22222) (0.76010,17.33333) (0.76335,17.44444) (0.76590,17.55556)
    (0.76773,17.66667) (0.76882,17.77778) (0.76916,17.88889) (0.76873,18.00000)
    (0.76752,18.11111) (0.76552,18.22222) (0.76273,18.33333) (0.75914,18.44444)
    (0.75475,18.55556) (0.74957,18.66667) (0.74361,18.77778) (0.73687,18.88889)
    (0.72938,19.00000) (0.72114,19.11111) (0.71219,19.22222) (0.70254,19.33333)
    (0.69223,19.44444) (0.68129,19.55556) (0.66975,19.66667) (0.65766,19.77778)
    (0.64506,19.88889) (0.63200,20.00000) (0.61853,20.11111) (0.60470,20.22222)
    (0.59056,20.33333) (0.57618,20.44444) (0.56162,20.55556) (0.54695,20.66667)
    (0.53222,20.77778) (0.51751,20.88889) (0.50288,21.00000) (0.50000,21.02204)
    (0.48842,21.11111) (0.47417,21.22222) (0.46023,21.33333) (0.44666,21.44444)
    (0.43353,21.55556) (0.42092,21.66667) (0.40889,21.77778) (0.39751,21.88889)
    (0.38685,22.00000) (0.37697,22.11111) (0.36794,22.22222) (0.35980,22.33333)
    (0.35263,22.44444) (0.34646,22.55556) (0.34134,22.66667) (0.33733,22.77778)
    (0.33444,22.88889) (0.33272,23.00000) (0.33219,23.11111) (0.33286,23.22222)
    (0.33477,23.33333) (0.33790,23.44444) (0.34226,23.55556) (0.34784,23.66667)
    (0.35464,23.77778) (0.36262,23.88889) (0.37176,24.00000) (0.38204,24.11111)
    (0.39339,24.22222) (0.40578,24.33333) (0.41915,24.44444) (0.43344,24.55556)
    (0.44858,24.66667) (0.46448,24.77778) (0.48108,24.88889) (0.49829,25.00000)
    (0.50000,25.01086) (0.51601,25.11111) (0.53415,25.22222) (0.55260,25.33333)
    (0.57127,25.44444) (0.59004,25.55556) (0.60882,25.66667) (0.62748,25.77778)
    (0.64591,25.88889) (0.66402,26.00000) (0.68167,26.11111) (0.69877,26.22222)
    (0.71520,26.33333) (0.73087,26.44444) (0.74566,26.55556) (0.75948,26.66667)
    (0.77224,26.77778) (0.78385,26.88889) (0.79423,27.00000) (0.80331,27.11111)
    (0.81102,27.22222) (0.81730,27.33333) (0.82211,27.44444) (0.82540,27.55556)
    (0.82716,27.66667) (0.82735,27.77778) (0.82597,27.88889) (0.82303,28.00000)
    (0.81853,28.11111) (0.81251,28.22222) (0.80499,28.33333) (0.79603,28.44444)
    (0.78568,28.55556) (0.77401,28.66667) (0.76110,28.77778) (0.74704,28.88889)
    (0.73193,29.00000) (0.71587,29.11111) (0.69898,29.22222) (0.68138,29.33333)
    (0.66320,29.44444) (0.64458,29.55556) (0.62566,29.66667) (0.60659,29.77778)
    (0.58750,29.88889) (0.56854,30.00000) (0.54988,30.11111) (0.53165,30.22222)
    (0.51400,30.33333) (0.50000,30.42488) (0.49708,30.44444) (0.48102,30.55556)
    (0.46597,30.66667) (0.45204,30.77778) (0.43937,30.88889) (0.42805,31.00000)
    (0.41820,31.11111) (0.40990,31.22222) (0.40323,31.33333) (0.39827,31.44444)
    (0.39506,31.55556) (0.39364,31.66667) (0.39404,31.77778) (0.39627,31.88889)
    (0.40032,32.00000) (0.40617,32.11111) (0.41380,32.22222) (0.42314,32.33333)
    (0.43413,32.44444) (0.44669,32.55556) (0.46072,32.66667) (0.47612,32.77778)
    (0.49275,32.88889) (0.50000,32.93506) (0.51050,33.00000) (0.52920,33.11111)
    (0.54870,33.22222) (0.56885,33.33333) (0.58945,33.44444) (0.61035,33.55556)
    (0.63135,33.66667) (0.65226,33.77778) (0.67290,33.88889) (0.69308,34.00000)
    (0.71261,34.11111) (0.73130,34.22222) (0.74899,34.33333) (0.76550,34.44444)
    (0.78067,34.55556) (0.79434,34.66667) (0.80638,34.77778) (0.81665,34.88889)
    (0.82505,35.00000) (0.83146,35.11111) (0.83582,35.22222) (0.83806,35.33333)
    (0.83814,35.44444) (0.83602,35.55556) (0.83171,35.66667) (0.82521,35.77778)
    (0.81655,35.88889) (0.80581,36.00000) (0.79303,36.11111) (0.77833,36.22222)
    (0.76180,36.33333) (0.74359,36.44444) (0.72382,36.55556) (0.70268,36.66667)
    (0.68032,36.77778) (0.65694,36.88889) (0.63274,37.00000) (0.60792,37.11111)
    (0.58270,37.22222) (0.55729,37.33333) (0.53193,37.44444) (0.50684,37.55556)
    (0.50000,37.58618) (0.48223,37.66667) (0.45833,37.77778) (0.43535,37.88889)
    (0.41350,38.00000) (0.39296,38.11111) (0.37392,38.22222) (0.35655,38.33333)
    (0.34100,38.44444) (0.32740,38.55556) (0.31587,38.66667) (0.30650,38.77778)
    (0.29936,38.88889) (0.29450,39.00000) (0.29195,39.11111) (0.29170,39.22222)
    (0.29374,39.33333) (0.29800,39.44444) (0.30443,39.55556) (0.31292,39.66667)
    (0.32336,39.77778) (0.33560,39.88889) (0.34948,40.00000)
  } -- (0.5,40.0) -- cycle;%
}
\providecommand{\CoverDrawHardyZCurve}[1][]{%
  \draw[#1] plot coordinates {
    (0.33206,0.00000) (0.33580,0.11111) (0.34557,0.22222) (0.35830,0.33333)
    (0.37134,0.44444) (0.38326,0.55556) (0.39355,0.66667) (0.40220,0.77778)
    (0.40939,0.88889) (0.41532,1.00000) (0.42022,1.11111) (0.42426,1.22222)
    (0.42759,1.33333) (0.43033,1.44444) (0.43258,1.55556) (0.43440,1.66667)
    (0.43587,1.77778) (0.43704,1.88889) (0.43794,2.00000) (0.43861,2.11111)
    (0.43908,2.22222) (0.43936,2.33333) (0.43948,2.44444) (0.43944,2.55556)
    (0.43928,2.66667) (0.43899,2.77778) (0.43858,2.88889) (0.43807,3.00000)
    (0.43745,3.11111) (0.43674,3.22222) (0.43594,3.33333) (0.43505,3.44444)
    (0.43408,3.55556) (0.43304,3.66667) (0.43191,3.77778) (0.43071,3.88889)
    (0.42944,4.00000) (0.42810,4.11111) (0.42669,4.22222) (0.42521,4.33333)
    (0.42367,4.44444) (0.42206,4.55556) (0.42039,4.66667) (0.41866,4.77778)
    (0.41688,4.88889) (0.41503,5.00000) (0.41313,5.11111) (0.41117,5.22222)
    (0.40916,5.33333) (0.40710,5.44444) (0.40499,5.55556) (0.40282,5.66667)
    (0.40062,5.77778) (0.39837,5.88889) (0.39607,6.00000) (0.39374,6.11111)
    (0.39137,6.22222) (0.38897,6.33333) (0.38653,6.44444) (0.38407,6.55556)
    (0.38158,6.66667) (0.37907,6.77778) (0.37655,6.88889) (0.37401,7.00000)
    (0.37146,7.11111) (0.36891,7.22222) (0.36635,7.33333) (0.36380,7.44444)
    (0.36127,7.55556) (0.35874,7.66667) (0.35624,7.77778) (0.35377,7.88889)
    (0.35133,8.00000) (0.34893,8.11111) (0.34658,8.22222) (0.34428,8.33333)
    (0.34205,8.44444) (0.33988,8.55556) (0.33779,8.66667) (0.33579,8.77778)
    (0.33388,8.88889) (0.33207,9.00000) (0.33037,9.11111) (0.32879,9.22222)
    (0.32734,9.33333) (0.32602,9.44444) (0.32485,9.55556) (0.32384,9.66667)
    (0.32300,9.77778) (0.32233,9.88889) (0.32184,10.00000) (0.32155,10.11111)
    (0.32147,10.22222) (0.32159,10.33333) (0.32195,10.44444) (0.32253,10.55556)
    (0.32336,10.66667) (0.32443,10.77778) (0.32577,10.88889) (0.32738,11.00000)
    (0.32926,11.11111) (0.33142,11.22222) (0.33388,11.33333) (0.33664,11.44444)
    (0.33970,11.55556) (0.34308,11.66667) (0.34677,11.77778) (0.35078,11.88889)
    (0.35511,12.00000) (0.35978,12.11111) (0.36477,12.22222) (0.37010,12.33333)
    (0.37576,12.44444) (0.38176,12.55556) (0.38808,12.66667) (0.39474,12.77778)
    (0.40172,12.88889) (0.40902,13.00000) (0.41664,13.11111) (0.42456,13.22222)
    (0.43279,13.33333) (0.44131,13.44444) (0.45011,13.55556) (0.45918,13.66667)
    (0.46850,13.77778) (0.47806,13.88889) (0.48785,14.00000) (0.49785,14.11111)
    (0.50000,14.13473) (0.50804,14.22222) (0.51839,14.33333) (0.52890,14.44444)
    (0.53952,14.55556) (0.55025,14.66667) (0.56106,14.77778) (0.57192,14.88889)
    (0.58279,15.00000) (0.59367,15.11111) (0.60450,15.22222) (0.61527,15.33333)
    (0.62595,15.44444) (0.63649,15.55556) (0.64687,15.66667) (0.65705,15.77778)
    (0.66701,15.88889) (0.67670,16.00000) (0.68610,16.11111) (0.69516,16.22222)
    (0.70385,16.33333) (0.71214,16.44444) (0.72000,16.55556) (0.72738,16.66667)
    (0.73427,16.77778) (0.74062,16.88889) (0.74641,17.00000) (0.75161,17.11111)
    (0.75618,17.22222) (0.76010,17.33333) (0.76335,17.44444) (0.76590,17.55556)
    (0.76773,17.66667) (0.76882,17.77778) (0.76916,17.88889) (0.76873,18.00000)
    (0.76752,18.11111) (0.76552,18.22222) (0.76273,18.33333) (0.75914,18.44444)
    (0.75475,18.55556) (0.74957,18.66667) (0.74361,18.77778) (0.73687,18.88889)
    (0.72938,19.00000) (0.72114,19.11111) (0.71219,19.22222) (0.70254,19.33333)
    (0.69223,19.44444) (0.68129,19.55556) (0.66975,19.66667) (0.65766,19.77778)
    (0.64506,19.88889) (0.63200,20.00000) (0.61853,20.11111) (0.60470,20.22222)
    (0.59056,20.33333) (0.57618,20.44444) (0.56162,20.55556) (0.54695,20.66667)
    (0.53222,20.77778) (0.51751,20.88889) (0.50288,21.00000) (0.50000,21.02204)
    (0.48842,21.11111) (0.47417,21.22222) (0.46023,21.33333) (0.44666,21.44444)
    (0.43353,21.55556) (0.42092,21.66667) (0.40889,21.77778) (0.39751,21.88889)
    (0.38685,22.00000) (0.37697,22.11111) (0.36794,22.22222) (0.35980,22.33333)
    (0.35263,22.44444) (0.34646,22.55556) (0.34134,22.66667) (0.33733,22.77778)
    (0.33444,22.88889) (0.33272,23.00000) (0.33219,23.11111) (0.33286,23.22222)
    (0.33477,23.33333) (0.33790,23.44444) (0.34226,23.55556) (0.34784,23.66667)
    (0.35464,23.77778) (0.36262,23.88889) (0.37176,24.00000) (0.38204,24.11111)
    (0.39339,24.22222) (0.40578,24.33333) (0.41915,24.44444) (0.43344,24.55556)
    (0.44858,24.66667) (0.46448,24.77778) (0.48108,24.88889) (0.49829,25.00000)
    (0.50000,25.01086) (0.51601,25.11111) (0.53415,25.22222) (0.55260,25.33333)
    (0.57127,25.44444) (0.59004,25.55556) (0.60882,25.66667) (0.62748,25.77778)
    (0.64591,25.88889) (0.66402,26.00000) (0.68167,26.11111) (0.69877,26.22222)
    (0.71520,26.33333) (0.73087,26.44444) (0.74566,26.55556) (0.75948,26.66667)
    (0.77224,26.77778) (0.78385,26.88889) (0.79423,27.00000) (0.80331,27.11111)
    (0.81102,27.22222) (0.81730,27.33333) (0.82211,27.44444) (0.82540,27.55556)
    (0.82716,27.66667) (0.82735,27.77778) (0.82597,27.88889) (0.82303,28.00000)
    (0.81853,28.11111) (0.81251,28.22222) (0.80499,28.33333) (0.79603,28.44444)
    (0.78568,28.55556) (0.77401,28.66667) (0.76110,28.77778) (0.74704,28.88889)
    (0.73193,29.00000) (0.71587,29.11111) (0.69898,29.22222) (0.68138,29.33333)
    (0.66320,29.44444) (0.64458,29.55556) (0.62566,29.66667) (0.60659,29.77778)
    (0.58750,29.88889) (0.56854,30.00000) (0.54988,30.11111) (0.53165,30.22222)
    (0.51400,30.33333) (0.50000,30.42488) (0.49708,30.44444) (0.48102,30.55556)
    (0.46597,30.66667) (0.45204,30.77778) (0.43937,30.88889) (0.42805,31.00000)
    (0.41820,31.11111) (0.40990,31.22222) (0.40323,31.33333) (0.39827,31.44444)
    (0.39506,31.55556) (0.39364,31.66667) (0.39404,31.77778) (0.39627,31.88889)
    (0.40032,32.00000) (0.40617,32.11111) (0.41380,32.22222) (0.42314,32.33333)
    (0.43413,32.44444) (0.44669,32.55556) (0.46072,32.66667) (0.47612,32.77778)
    (0.49275,32.88889) (0.50000,32.93506) (0.51050,33.00000) (0.52920,33.11111)
    (0.54870,33.22222) (0.56885,33.33333) (0.58945,33.44444) (0.61035,33.55556)
    (0.63135,33.66667) (0.65226,33.77778) (0.67290,33.88889) (0.69308,34.00000)
    (0.71261,34.11111) (0.73130,34.22222) (0.74899,34.33333) (0.76550,34.44444)
    (0.78067,34.55556) (0.79434,34.66667) (0.80638,34.77778) (0.81665,34.88889)
    (0.82505,35.00000) (0.83146,35.11111) (0.83582,35.22222) (0.83806,35.33333)
    (0.83814,35.44444) (0.83602,35.55556) (0.83171,35.66667) (0.82521,35.77778)
    (0.81655,35.88889) (0.80581,36.00000) (0.79303,36.11111) (0.77833,36.22222)
    (0.76180,36.33333) (0.74359,36.44444) (0.72382,36.55556) (0.70268,36.66667)
    (0.68032,36.77778) (0.65694,36.88889) (0.63274,37.00000) (0.60792,37.11111)
    (0.58270,37.22222) (0.55729,37.33333) (0.53193,37.44444) (0.50684,37.55556)
    (0.50000,37.58618) (0.48223,37.66667) (0.45833,37.77778) (0.43535,37.88889)
    (0.41350,38.00000) (0.39296,38.11111) (0.37392,38.22222) (0.35655,38.33333)
    (0.34100,38.44444) (0.32740,38.55556) (0.31587,38.66667) (0.30650,38.77778)
    (0.29936,38.88889) (0.29450,39.00000) (0.29195,39.11111) (0.29170,39.22222)
    (0.29374,39.33333) (0.29800,39.44444) (0.30443,39.55556) (0.31292,39.66667)
    (0.32336,39.77778) (0.33560,39.88889) (0.34948,40.00000)
  };%
}
%
}{%
  % Auto-generated by plot_critical_strip.py using 50-digit mpmath arithmetic.
% Coordinates are in critical-strip units: x = Re(s), y = Im(s) on [0, 40].
\providecommand{\CoverCriticalZeros}{%
  14.134725, 21.022040, 25.010858, 30.424876, 32.935062, 37.586178%
}
\providecommand{\CoverFillHardyZArea}[1][]{%
  \fill[#1] (0.5,0) -- plot coordinates {
    (0.33206,0.00000) (0.33580,0.11111) (0.34557,0.22222) (0.35830,0.33333)
    (0.37134,0.44444) (0.38326,0.55556) (0.39355,0.66667) (0.40220,0.77778)
    (0.40939,0.88889) (0.41532,1.00000) (0.42022,1.11111) (0.42426,1.22222)
    (0.42759,1.33333) (0.43033,1.44444) (0.43258,1.55556) (0.43440,1.66667)
    (0.43587,1.77778) (0.43704,1.88889) (0.43794,2.00000) (0.43861,2.11111)
    (0.43908,2.22222) (0.43936,2.33333) (0.43948,2.44444) (0.43944,2.55556)
    (0.43928,2.66667) (0.43899,2.77778) (0.43858,2.88889) (0.43807,3.00000)
    (0.43745,3.11111) (0.43674,3.22222) (0.43594,3.33333) (0.43505,3.44444)
    (0.43408,3.55556) (0.43304,3.66667) (0.43191,3.77778) (0.43071,3.88889)
    (0.42944,4.00000) (0.42810,4.11111) (0.42669,4.22222) (0.42521,4.33333)
    (0.42367,4.44444) (0.42206,4.55556) (0.42039,4.66667) (0.41866,4.77778)
    (0.41688,4.88889) (0.41503,5.00000) (0.41313,5.11111) (0.41117,5.22222)
    (0.40916,5.33333) (0.40710,5.44444) (0.40499,5.55556) (0.40282,5.66667)
    (0.40062,5.77778) (0.39837,5.88889) (0.39607,6.00000) (0.39374,6.11111)
    (0.39137,6.22222) (0.38897,6.33333) (0.38653,6.44444) (0.38407,6.55556)
    (0.38158,6.66667) (0.37907,6.77778) (0.37655,6.88889) (0.37401,7.00000)
    (0.37146,7.11111) (0.36891,7.22222) (0.36635,7.33333) (0.36380,7.44444)
    (0.36127,7.55556) (0.35874,7.66667) (0.35624,7.77778) (0.35377,7.88889)
    (0.35133,8.00000) (0.34893,8.11111) (0.34658,8.22222) (0.34428,8.33333)
    (0.34205,8.44444) (0.33988,8.55556) (0.33779,8.66667) (0.33579,8.77778)
    (0.33388,8.88889) (0.33207,9.00000) (0.33037,9.11111) (0.32879,9.22222)
    (0.32734,9.33333) (0.32602,9.44444) (0.32485,9.55556) (0.32384,9.66667)
    (0.32300,9.77778) (0.32233,9.88889) (0.32184,10.00000) (0.32155,10.11111)
    (0.32147,10.22222) (0.32159,10.33333) (0.32195,10.44444) (0.32253,10.55556)
    (0.32336,10.66667) (0.32443,10.77778) (0.32577,10.88889) (0.32738,11.00000)
    (0.32926,11.11111) (0.33142,11.22222) (0.33388,11.33333) (0.33664,11.44444)
    (0.33970,11.55556) (0.34308,11.66667) (0.34677,11.77778) (0.35078,11.88889)
    (0.35511,12.00000) (0.35978,12.11111) (0.36477,12.22222) (0.37010,12.33333)
    (0.37576,12.44444) (0.38176,12.55556) (0.38808,12.66667) (0.39474,12.77778)
    (0.40172,12.88889) (0.40902,13.00000) (0.41664,13.11111) (0.42456,13.22222)
    (0.43279,13.33333) (0.44131,13.44444) (0.45011,13.55556) (0.45918,13.66667)
    (0.46850,13.77778) (0.47806,13.88889) (0.48785,14.00000) (0.49785,14.11111)
    (0.50000,14.13473) (0.50804,14.22222) (0.51839,14.33333) (0.52890,14.44444)
    (0.53952,14.55556) (0.55025,14.66667) (0.56106,14.77778) (0.57192,14.88889)
    (0.58279,15.00000) (0.59367,15.11111) (0.60450,15.22222) (0.61527,15.33333)
    (0.62595,15.44444) (0.63649,15.55556) (0.64687,15.66667) (0.65705,15.77778)
    (0.66701,15.88889) (0.67670,16.00000) (0.68610,16.11111) (0.69516,16.22222)
    (0.70385,16.33333) (0.71214,16.44444) (0.72000,16.55556) (0.72738,16.66667)
    (0.73427,16.77778) (0.74062,16.88889) (0.74641,17.00000) (0.75161,17.11111)
    (0.75618,17.22222) (0.76010,17.33333) (0.76335,17.44444) (0.76590,17.55556)
    (0.76773,17.66667) (0.76882,17.77778) (0.76916,17.88889) (0.76873,18.00000)
    (0.76752,18.11111) (0.76552,18.22222) (0.76273,18.33333) (0.75914,18.44444)
    (0.75475,18.55556) (0.74957,18.66667) (0.74361,18.77778) (0.73687,18.88889)
    (0.72938,19.00000) (0.72114,19.11111) (0.71219,19.22222) (0.70254,19.33333)
    (0.69223,19.44444) (0.68129,19.55556) (0.66975,19.66667) (0.65766,19.77778)
    (0.64506,19.88889) (0.63200,20.00000) (0.61853,20.11111) (0.60470,20.22222)
    (0.59056,20.33333) (0.57618,20.44444) (0.56162,20.55556) (0.54695,20.66667)
    (0.53222,20.77778) (0.51751,20.88889) (0.50288,21.00000) (0.50000,21.02204)
    (0.48842,21.11111) (0.47417,21.22222) (0.46023,21.33333) (0.44666,21.44444)
    (0.43353,21.55556) (0.42092,21.66667) (0.40889,21.77778) (0.39751,21.88889)
    (0.38685,22.00000) (0.37697,22.11111) (0.36794,22.22222) (0.35980,22.33333)
    (0.35263,22.44444) (0.34646,22.55556) (0.34134,22.66667) (0.33733,22.77778)
    (0.33444,22.88889) (0.33272,23.00000) (0.33219,23.11111) (0.33286,23.22222)
    (0.33477,23.33333) (0.33790,23.44444) (0.34226,23.55556) (0.34784,23.66667)
    (0.35464,23.77778) (0.36262,23.88889) (0.37176,24.00000) (0.38204,24.11111)
    (0.39339,24.22222) (0.40578,24.33333) (0.41915,24.44444) (0.43344,24.55556)
    (0.44858,24.66667) (0.46448,24.77778) (0.48108,24.88889) (0.49829,25.00000)
    (0.50000,25.01086) (0.51601,25.11111) (0.53415,25.22222) (0.55260,25.33333)
    (0.57127,25.44444) (0.59004,25.55556) (0.60882,25.66667) (0.62748,25.77778)
    (0.64591,25.88889) (0.66402,26.00000) (0.68167,26.11111) (0.69877,26.22222)
    (0.71520,26.33333) (0.73087,26.44444) (0.74566,26.55556) (0.75948,26.66667)
    (0.77224,26.77778) (0.78385,26.88889) (0.79423,27.00000) (0.80331,27.11111)
    (0.81102,27.22222) (0.81730,27.33333) (0.82211,27.44444) (0.82540,27.55556)
    (0.82716,27.66667) (0.82735,27.77778) (0.82597,27.88889) (0.82303,28.00000)
    (0.81853,28.11111) (0.81251,28.22222) (0.80499,28.33333) (0.79603,28.44444)
    (0.78568,28.55556) (0.77401,28.66667) (0.76110,28.77778) (0.74704,28.88889)
    (0.73193,29.00000) (0.71587,29.11111) (0.69898,29.22222) (0.68138,29.33333)
    (0.66320,29.44444) (0.64458,29.55556) (0.62566,29.66667) (0.60659,29.77778)
    (0.58750,29.88889) (0.56854,30.00000) (0.54988,30.11111) (0.53165,30.22222)
    (0.51400,30.33333) (0.50000,30.42488) (0.49708,30.44444) (0.48102,30.55556)
    (0.46597,30.66667) (0.45204,30.77778) (0.43937,30.88889) (0.42805,31.00000)
    (0.41820,31.11111) (0.40990,31.22222) (0.40323,31.33333) (0.39827,31.44444)
    (0.39506,31.55556) (0.39364,31.66667) (0.39404,31.77778) (0.39627,31.88889)
    (0.40032,32.00000) (0.40617,32.11111) (0.41380,32.22222) (0.42314,32.33333)
    (0.43413,32.44444) (0.44669,32.55556) (0.46072,32.66667) (0.47612,32.77778)
    (0.49275,32.88889) (0.50000,32.93506) (0.51050,33.00000) (0.52920,33.11111)
    (0.54870,33.22222) (0.56885,33.33333) (0.58945,33.44444) (0.61035,33.55556)
    (0.63135,33.66667) (0.65226,33.77778) (0.67290,33.88889) (0.69308,34.00000)
    (0.71261,34.11111) (0.73130,34.22222) (0.74899,34.33333) (0.76550,34.44444)
    (0.78067,34.55556) (0.79434,34.66667) (0.80638,34.77778) (0.81665,34.88889)
    (0.82505,35.00000) (0.83146,35.11111) (0.83582,35.22222) (0.83806,35.33333)
    (0.83814,35.44444) (0.83602,35.55556) (0.83171,35.66667) (0.82521,35.77778)
    (0.81655,35.88889) (0.80581,36.00000) (0.79303,36.11111) (0.77833,36.22222)
    (0.76180,36.33333) (0.74359,36.44444) (0.72382,36.55556) (0.70268,36.66667)
    (0.68032,36.77778) (0.65694,36.88889) (0.63274,37.00000) (0.60792,37.11111)
    (0.58270,37.22222) (0.55729,37.33333) (0.53193,37.44444) (0.50684,37.55556)
    (0.50000,37.58618) (0.48223,37.66667) (0.45833,37.77778) (0.43535,37.88889)
    (0.41350,38.00000) (0.39296,38.11111) (0.37392,38.22222) (0.35655,38.33333)
    (0.34100,38.44444) (0.32740,38.55556) (0.31587,38.66667) (0.30650,38.77778)
    (0.29936,38.88889) (0.29450,39.00000) (0.29195,39.11111) (0.29170,39.22222)
    (0.29374,39.33333) (0.29800,39.44444) (0.30443,39.55556) (0.31292,39.66667)
    (0.32336,39.77778) (0.33560,39.88889) (0.34948,40.00000)
  } -- (0.5,40.0) -- cycle;%
}
\providecommand{\CoverDrawHardyZCurve}[1][]{%
  \draw[#1] plot coordinates {
    (0.33206,0.00000) (0.33580,0.11111) (0.34557,0.22222) (0.35830,0.33333)
    (0.37134,0.44444) (0.38326,0.55556) (0.39355,0.66667) (0.40220,0.77778)
    (0.40939,0.88889) (0.41532,1.00000) (0.42022,1.11111) (0.42426,1.22222)
    (0.42759,1.33333) (0.43033,1.44444) (0.43258,1.55556) (0.43440,1.66667)
    (0.43587,1.77778) (0.43704,1.88889) (0.43794,2.00000) (0.43861,2.11111)
    (0.43908,2.22222) (0.43936,2.33333) (0.43948,2.44444) (0.43944,2.55556)
    (0.43928,2.66667) (0.43899,2.77778) (0.43858,2.88889) (0.43807,3.00000)
    (0.43745,3.11111) (0.43674,3.22222) (0.43594,3.33333) (0.43505,3.44444)
    (0.43408,3.55556) (0.43304,3.66667) (0.43191,3.77778) (0.43071,3.88889)
    (0.42944,4.00000) (0.42810,4.11111) (0.42669,4.22222) (0.42521,4.33333)
    (0.42367,4.44444) (0.42206,4.55556) (0.42039,4.66667) (0.41866,4.77778)
    (0.41688,4.88889) (0.41503,5.00000) (0.41313,5.11111) (0.41117,5.22222)
    (0.40916,5.33333) (0.40710,5.44444) (0.40499,5.55556) (0.40282,5.66667)
    (0.40062,5.77778) (0.39837,5.88889) (0.39607,6.00000) (0.39374,6.11111)
    (0.39137,6.22222) (0.38897,6.33333) (0.38653,6.44444) (0.38407,6.55556)
    (0.38158,6.66667) (0.37907,6.77778) (0.37655,6.88889) (0.37401,7.00000)
    (0.37146,7.11111) (0.36891,7.22222) (0.36635,7.33333) (0.36380,7.44444)
    (0.36127,7.55556) (0.35874,7.66667) (0.35624,7.77778) (0.35377,7.88889)
    (0.35133,8.00000) (0.34893,8.11111) (0.34658,8.22222) (0.34428,8.33333)
    (0.34205,8.44444) (0.33988,8.55556) (0.33779,8.66667) (0.33579,8.77778)
    (0.33388,8.88889) (0.33207,9.00000) (0.33037,9.11111) (0.32879,9.22222)
    (0.32734,9.33333) (0.32602,9.44444) (0.32485,9.55556) (0.32384,9.66667)
    (0.32300,9.77778) (0.32233,9.88889) (0.32184,10.00000) (0.32155,10.11111)
    (0.32147,10.22222) (0.32159,10.33333) (0.32195,10.44444) (0.32253,10.55556)
    (0.32336,10.66667) (0.32443,10.77778) (0.32577,10.88889) (0.32738,11.00000)
    (0.32926,11.11111) (0.33142,11.22222) (0.33388,11.33333) (0.33664,11.44444)
    (0.33970,11.55556) (0.34308,11.66667) (0.34677,11.77778) (0.35078,11.88889)
    (0.35511,12.00000) (0.35978,12.11111) (0.36477,12.22222) (0.37010,12.33333)
    (0.37576,12.44444) (0.38176,12.55556) (0.38808,12.66667) (0.39474,12.77778)
    (0.40172,12.88889) (0.40902,13.00000) (0.41664,13.11111) (0.42456,13.22222)
    (0.43279,13.33333) (0.44131,13.44444) (0.45011,13.55556) (0.45918,13.66667)
    (0.46850,13.77778) (0.47806,13.88889) (0.48785,14.00000) (0.49785,14.11111)
    (0.50000,14.13473) (0.50804,14.22222) (0.51839,14.33333) (0.52890,14.44444)
    (0.53952,14.55556) (0.55025,14.66667) (0.56106,14.77778) (0.57192,14.88889)
    (0.58279,15.00000) (0.59367,15.11111) (0.60450,15.22222) (0.61527,15.33333)
    (0.62595,15.44444) (0.63649,15.55556) (0.64687,15.66667) (0.65705,15.77778)
    (0.66701,15.88889) (0.67670,16.00000) (0.68610,16.11111) (0.69516,16.22222)
    (0.70385,16.33333) (0.71214,16.44444) (0.72000,16.55556) (0.72738,16.66667)
    (0.73427,16.77778) (0.74062,16.88889) (0.74641,17.00000) (0.75161,17.11111)
    (0.75618,17.22222) (0.76010,17.33333) (0.76335,17.44444) (0.76590,17.55556)
    (0.76773,17.66667) (0.76882,17.77778) (0.76916,17.88889) (0.76873,18.00000)
    (0.76752,18.11111) (0.76552,18.22222) (0.76273,18.33333) (0.75914,18.44444)
    (0.75475,18.55556) (0.74957,18.66667) (0.74361,18.77778) (0.73687,18.88889)
    (0.72938,19.00000) (0.72114,19.11111) (0.71219,19.22222) (0.70254,19.33333)
    (0.69223,19.44444) (0.68129,19.55556) (0.66975,19.66667) (0.65766,19.77778)
    (0.64506,19.88889) (0.63200,20.00000) (0.61853,20.11111) (0.60470,20.22222)
    (0.59056,20.33333) (0.57618,20.44444) (0.56162,20.55556) (0.54695,20.66667)
    (0.53222,20.77778) (0.51751,20.88889) (0.50288,21.00000) (0.50000,21.02204)
    (0.48842,21.11111) (0.47417,21.22222) (0.46023,21.33333) (0.44666,21.44444)
    (0.43353,21.55556) (0.42092,21.66667) (0.40889,21.77778) (0.39751,21.88889)
    (0.38685,22.00000) (0.37697,22.11111) (0.36794,22.22222) (0.35980,22.33333)
    (0.35263,22.44444) (0.34646,22.55556) (0.34134,22.66667) (0.33733,22.77778)
    (0.33444,22.88889) (0.33272,23.00000) (0.33219,23.11111) (0.33286,23.22222)
    (0.33477,23.33333) (0.33790,23.44444) (0.34226,23.55556) (0.34784,23.66667)
    (0.35464,23.77778) (0.36262,23.88889) (0.37176,24.00000) (0.38204,24.11111)
    (0.39339,24.22222) (0.40578,24.33333) (0.41915,24.44444) (0.43344,24.55556)
    (0.44858,24.66667) (0.46448,24.77778) (0.48108,24.88889) (0.49829,25.00000)
    (0.50000,25.01086) (0.51601,25.11111) (0.53415,25.22222) (0.55260,25.33333)
    (0.57127,25.44444) (0.59004,25.55556) (0.60882,25.66667) (0.62748,25.77778)
    (0.64591,25.88889) (0.66402,26.00000) (0.68167,26.11111) (0.69877,26.22222)
    (0.71520,26.33333) (0.73087,26.44444) (0.74566,26.55556) (0.75948,26.66667)
    (0.77224,26.77778) (0.78385,26.88889) (0.79423,27.00000) (0.80331,27.11111)
    (0.81102,27.22222) (0.81730,27.33333) (0.82211,27.44444) (0.82540,27.55556)
    (0.82716,27.66667) (0.82735,27.77778) (0.82597,27.88889) (0.82303,28.00000)
    (0.81853,28.11111) (0.81251,28.22222) (0.80499,28.33333) (0.79603,28.44444)
    (0.78568,28.55556) (0.77401,28.66667) (0.76110,28.77778) (0.74704,28.88889)
    (0.73193,29.00000) (0.71587,29.11111) (0.69898,29.22222) (0.68138,29.33333)
    (0.66320,29.44444) (0.64458,29.55556) (0.62566,29.66667) (0.60659,29.77778)
    (0.58750,29.88889) (0.56854,30.00000) (0.54988,30.11111) (0.53165,30.22222)
    (0.51400,30.33333) (0.50000,30.42488) (0.49708,30.44444) (0.48102,30.55556)
    (0.46597,30.66667) (0.45204,30.77778) (0.43937,30.88889) (0.42805,31.00000)
    (0.41820,31.11111) (0.40990,31.22222) (0.40323,31.33333) (0.39827,31.44444)
    (0.39506,31.55556) (0.39364,31.66667) (0.39404,31.77778) (0.39627,31.88889)
    (0.40032,32.00000) (0.40617,32.11111) (0.41380,32.22222) (0.42314,32.33333)
    (0.43413,32.44444) (0.44669,32.55556) (0.46072,32.66667) (0.47612,32.77778)
    (0.49275,32.88889) (0.50000,32.93506) (0.51050,33.00000) (0.52920,33.11111)
    (0.54870,33.22222) (0.56885,33.33333) (0.58945,33.44444) (0.61035,33.55556)
    (0.63135,33.66667) (0.65226,33.77778) (0.67290,33.88889) (0.69308,34.00000)
    (0.71261,34.11111) (0.73130,34.22222) (0.74899,34.33333) (0.76550,34.44444)
    (0.78067,34.55556) (0.79434,34.66667) (0.80638,34.77778) (0.81665,34.88889)
    (0.82505,35.00000) (0.83146,35.11111) (0.83582,35.22222) (0.83806,35.33333)
    (0.83814,35.44444) (0.83602,35.55556) (0.83171,35.66667) (0.82521,35.77778)
    (0.81655,35.88889) (0.80581,36.00000) (0.79303,36.11111) (0.77833,36.22222)
    (0.76180,36.33333) (0.74359,36.44444) (0.72382,36.55556) (0.70268,36.66667)
    (0.68032,36.77778) (0.65694,36.88889) (0.63274,37.00000) (0.60792,37.11111)
    (0.58270,37.22222) (0.55729,37.33333) (0.53193,37.44444) (0.50684,37.55556)
    (0.50000,37.58618) (0.48223,37.66667) (0.45833,37.77778) (0.43535,37.88889)
    (0.41350,38.00000) (0.39296,38.11111) (0.37392,38.22222) (0.35655,38.33333)
    (0.34100,38.44444) (0.32740,38.55556) (0.31587,38.66667) (0.30650,38.77778)
    (0.29936,38.88889) (0.29450,39.00000) (0.29195,39.11111) (0.29170,39.22222)
    (0.29374,39.33333) (0.29800,39.44444) (0.30443,39.55556) (0.31292,39.66667)
    (0.32336,39.77778) (0.33560,39.88889) (0.34948,40.00000)
  };%
}
%
}

\begin{tikzpicture}[
  x=1pt, y=1pt,
  font=\CoverSerifFont\fontsize{9.2}{11.2}\selectfont,
  text=CoverInk,
  eyebrow/.style={
    anchor=west, inner sep=0pt,
    font=\CoverSansFont\bfseries\addfontfeature{LetterSpace=7.0}\fontsize{7.6}{9.5}\selectfont,
    text=CoverInk
  },
  tier/.style={
    sharp corners, fill=white, draw=CoverHairline, line width=0.4pt,
    minimum width=178pt, minimum height=27pt, inner sep=0pt,
    font=\CoverSerifFont\fontsize{9.8}{12}\selectfont, text=CoverInk
  },
  flow arrow/.style={
    -{Stealth[length=4.0pt,width=2.8pt]}, draw=CoverAxis, line width=0.65pt
  },
  note/.style={
    inner sep=0pt, text=CoverMuted,
    font=\CoverSerifFont\fontsize{8.5}{10.8}\selectfont
  },
  step label/.style={
    anchor=north, align=center, inner sep=0pt,
    font=\CoverSerifFont\fontsize{8.5}{10.4}\selectfont, text=CoverInk
  },
  edge note/.style={
    anchor=south, inner sep=0pt,
    font=\CoverSerifFont\itshape\fontsize{8.0}{9.4}\selectfont, text=CoverMuted
  },
  tick label/.style={
    inner sep=0pt,
    font=\CoverSansFont\fontsize{7.5}{9}\selectfont, text=CoverMuted
  },
]
  % Full textwidth (435pt) x grand portrait-balanced height (346pt), borderless light-gray canvas
  \path[use as bounding box] (0,0) rectangle (435,346);
  \fill[CoverBg] (0,0) rectangle (435,346);

  % Quiet interior hairlines only (no heavy top/bottom border lines)
  \draw[CoverRule, line width=0.45pt] (16,108)  -- (419,108);
  \draw[CoverRule, line width=0.45pt] (214,122) -- (214,330);

  % ==========================================================================
  % 01 | THETA LIFTS (Top Left, x = 0..214, y = 108..346)
  % ==========================================================================
  \node[eyebrow] at (16,328) {%
    {\color{CoverLightInk}01}\hspace{5pt}{\color{CoverRule}\rule[-0.12ex]{0.55pt}{1.15ex}}\hspace{5pt}THETA LIFTS%
  };

  % Source tier (crisp white band with calm mid-gray left accent)
  \node[tier] (source) at (107,294) {Automorphic form \(f\) on \(G\)};
  \fill[CoverTierLight] (source.south west) rectangle ($(source.north west)+(1.8pt,0)$);

  % Sculptural symmetric pairing arcs framing the integral operator
  \draw[CoverRule, line width=0.48pt]
    (21,274) .. controls (13,242) and (13,198) .. (21,166);
  \draw[CoverRule, line width=0.48pt]
    (193,274) .. controls (201,242) and (201,198) .. (193,166);
  \fill[CoverRule] (21,274) circle[radius=1.1pt]  (21,166) circle[radius=1.1pt]
                   (193,274) circle[radius=1.1pt] (193,166) circle[radius=1.1pt];

  % Central integral kernel (generous symmetrical 24pt flow arrows)
  \draw[flow arrow] (107,277.5) -- (107,253.5);
  \node[font=\CoverSerifFont\fontsize{10.0}{12}\selectfont, text=CoverInk] (integral) at (107,231) {%
    \CoverLink{eq:adelic-theta-integral}{\(\displaystyle
      (\Theta_\varphi f)(h)=\int_{[G]} f(g)\,\theta_\varphi(g,h)\,dg\)}%
  };
  \node[note, align=center] at (107,194) {%
    Integrate against the theta kernel\\
    over the source quotient \([G]\)%
  };
  \draw[flow arrow] (107,178.5) -- (107,154.5);

  % Target tier (deeper slate-gray left accent showing lift progression)
  \node[tier] (target) at (107,138) {Lifted form \(\Theta_\varphi f\) on \(H\)};
  \fill[CoverTierDark] (target.south west) rectangle ($(target.north west)+(1.8pt,0)$);

  % ==========================================================================
  % 02 | THE CRITICAL LINE (Top Right, x = 214..435, y = 108..346)
  % ==========================================================================
  \node[eyebrow] at (230,328) {%
    {\color{CoverLightInk}02}\hspace{5pt}{\color{CoverRule}\rule[-0.12ex]{0.55pt}{1.15ex}}\hspace{5pt}THE CRITICAL LINE%
  };

  % Tall native TikZ critical-strip coordinate system:
  % Re(s) = 0 -> x = 264pt, Re(s) = 1/2 -> x = 324pt, Re(s) = 1 -> x = 384pt (1 unit = 120pt)
  % Im(s) = 0 -> y = 179pt, Im(s) = 40  -> y = 301pt (1 unit = 3.05pt, 122pt tall!)
  \begin{scope}[shift={(264,179)}, x=120pt, y=3.05pt]
    % Critical strip 0 <= Re(s) <= 1 (subtle tonal split at Re(s) = 1/2)
    \fill[CoverStripLeft]  (0,0)   rectangle (0.5,40.4);
    \fill[CoverStripRight] (0.5,0) rectangle (1,40.4);

    % Subtle horizontal grid lines
    \foreach \yt in {10,20,30,40} {
      \draw[CoverGrid, line width=0.38pt] (-0.15,\yt) -- (1.10,\yt);
    }

    % Strip boundaries Re(s) = 0 and Re(s) = 1
    \draw[CoverRule, line width=0.55pt] (0,0) -- (0,40.4);
    \draw[CoverRule, line width=0.55pt] (1,0) -- (1,40.4);

    % Exact Hardy Z(t) wave along the critical line (shaded + stroke)
    \CoverFillHardyZArea[fill=CoverWaveFill, fill opacity=0.14]
    \CoverDrawHardyZCurve[draw=CoverWaveLine, line width=0.82pt, line join=round]

    % Critical line Re(s) = 1/2 (sole restrained burgundy focal line, flush with strip top)
    \draw[DeepBurgundy, line width=1.0pt, dash pattern=on 2.8pt off 2.0pt]
      (0.5,0) -- (0.5,40.4);

    % Axes (L-shaped, crisp open gray)
    \draw[CoverAxis, line width=0.58pt, -{Stealth[length=3.4pt,width=2.3pt]}]
      (-0.15,0) -- (1.13,0);
    \draw[CoverAxis, line width=0.58pt, -{Stealth[length=3.4pt,width=2.3pt]}]
      (-0.15,0) -- (-0.15,41.6);

    % Y-axis ticks & labels
    \foreach \yt in {0,10,20,30,40} {
      \draw[CoverAxis, line width=0.52pt] (-0.15,\yt) -- ++(-2.4pt,0);
      \node[tick label, anchor=east] at ($(-0.15,\yt)+(-4.4pt,0)$) {\yt};
    }

    % X-axis ticks & labels
    \foreach \xt/\xl in {0/0, 1/1} {
      \draw[CoverAxis, line width=0.52pt] (\xt,0) -- ++(0,-2.4pt);
      \node[tick label, anchor=north] at ($(\xt,0)+(0,-4.4pt)$) {\xl};
    }
    \draw[DeepBurgundy, line width=0.68pt] (0.5,0) -- ++(0,-2.7pt);
    \node[tick label, anchor=north, text=DeepBurgundy,
          font=\CoverSansFont\bfseries\fontsize{7.7}{9}\selectfont]
      at ($(0.5,0)+(0,-4.4pt)$) {1/2};

    % Axis titles (matched TeX math font, generous clearance from eyebrow & captions)
    \node[anchor=south, inner sep=0pt, text=CoverInk,
          font=\CoverSerifFont\fontsize{8.1}{9.5}\selectfont]
      at ($(-0.15,41.6)+(0,2.4pt)$) {\(\operatorname{Im}(s)\)};
    \node[anchor=west, inner sep=0pt, text=CoverInk,
          font=\CoverSerifFont\fontsize{8.1}{9.5}\selectfont]
      at ($(1.13,0)+(2.4pt,0)$) {\(\operatorname{Re}(s)\)};

    % Subtle Hardy Z(t) annotation in calm lower-left strip space
    \node[anchor=east, inner sep=0pt, text=CoverLightInk,
          font=\CoverSerifFont\fontsize{7.5}{9}\selectfont]
      at (0.26,7.2) {\(Z(t)\)};

    % First six nontrivial zeta zeros on Re(s) = 1/2
    \foreach \zt in \CoverCriticalZeros {
      \filldraw[fill=DeepBurgundy, draw=white, line width=0.7pt]
        (0.5,\zt) circle[radius=2.35pt];
    }
  \end{scope}

  % Right-column captions (clear hierarchy & symmetric cushion above y = 108)
  \node[note] at (324,154.5) {First six nontrivial zeros with \(\operatorname{Im}(s)>0\)};
  \node[align=center, inner sep=0pt,
        font=\CoverSerifFont\fontsize{9.0}{11.4}\selectfont, text=CoverInk] at (324,132.5) {%
    \CoverLink{thm:rh}{\textcolor{DeepBurgundy}{\textbf{RH}}}. Every nontrivial zero of \(\zeta(s)\) lies\\
    on the critical line \(\operatorname{Re}(s)=1/2\).%
  };

  % ==========================================================================
  % 03 | FROM THETA LIFTS TO THE PROOF OF GRH AND RH (Bottom, x = 0..435, y = 0..108)
  % ==========================================================================
  \node[eyebrow] at (16,92) {%
    {\color{CoverLightInk}03}\hspace{5pt}{\color{CoverRule}\rule[-0.12ex]{0.55pt}{1.15ex}}\hspace{5pt}FROM THETA LIFTS TO THE PROOF OF GRH AND RH%
  };

  % Horizontal proof-roadmap timeline arrows (y = 64, uniform 72pt pitch)
  \draw[flow arrow] (46.0,64)  -- (103.8,64);
  \draw[flow arrow] (115.8,64) -- (175.8,64);
  \draw[flow arrow] (187.8,64) -- (247.8,64);
  \draw[flow arrow] (259.8,64) -- (319.8,64);
  \draw[flow arrow, draw=DeepBurgundy!65!CoverAxis] (331.8,64) -- (389.2,64);

  % Clean italic arrow annotations (ample separation, clear of Node 6 halo)
  \node[edge note] at (290,71.5) {Landau criterion};
  \node[edge note] at (360,71.5) {Functional equations};

  % Timeline nodes
  % 1. Starting engine (filled slate-gray + soft ring)
  \draw[CoverAxis!55, line width=0.58pt] (38,64) circle[radius=5.8pt];
  \fill[CoverTierDark] (38,64) circle[radius=3.5pt];

  % 2-5. Intermediate theorem stages (precision hollow ring + center dot)
  \foreach \nx in {110,182,254,326} {
    \filldraw[fill=white, draw=CoverAxis, line width=0.85pt] (\nx,64) circle[radius=3.5pt];
    \fill[CoverTierDark] (\nx,64) circle[radius=1.15pt];
  }

  % 6. Destination (DeepBurgundy double-ring accent)
  \draw[DeepBurgundy!65, line width=0.7pt] (398,64) circle[radius=6.6pt];
  \fill[DeepBurgundy] (398,64) circle[radius=3.9pt];

  % Step labels with all original PDF hyperlinks preserved (generous bottom margin)
  \node[step label] at (38,52) {%
    \CoverLink{sec:theta-lifts}{Iterated}\\
    \CoverLink{sec:theta-lifts}{theta lifts}%
  };
  \node[step label] at (110,52) {%
    \CoverLink{sec:main-inequality}{Theorems}\\
    \CoverLink{sec:main-inequality}{A–C}%
  };
  \node[step label] at (182,52) {%
    \CoverLink{thm:star}{Main}\\
    \CoverLink{thm:star}{inequality}%
  };
  \node[step label] at (254,52) {%
    \CoverLink{thm:one-sided}{One-sided}\\
    \CoverLink{thm:one-sided}{prime-power}\\
    \CoverLink{thm:one-sided}{bound}%
  };
  \node[step label] at (326,52) {%
    \CoverLink{thm:half-plane}{Zero-free}\\
    \CoverLink{thm:half-plane}{half-plane}\\
    \CoverLink{thm:half-plane}{\(\operatorname{Re}(s)>1/2\)}%
  };
  \node[step label] at (398,52) {%
    \CoverLink{thm:rh}{\textcolor{DeepBurgundy}{\textbf{RH}}} and\\
    \CoverLink{thm:l-zeros}{\textcolor{DeepBurgundy}{\textbf{Dirichlet GRH}}}%
  };
\end{tikzpicture}
\par}
\end{titlepage}


\chapter*{Abstract}
\addcontentsline{toc}{chapter}{Abstract}
\markboth{Abstract}{Abstract}

This monograph follows a streamlined Lean blueprint in giving a systematic account of the complete argument for the Riemann hypothesis and the generalized Riemann hypothesis for Dirichlet $L$-functions. The blueprint is based on the Lean~4 formalization by Levent Alpöge and Claude. Starting from admissible data, weighted prime-power remainders, and adelic and algebraic theta lifts on unitary groups, the book establishes four arithmetic and analytic estimates. \statementlink{thm:short-forms}{Theorem~A(a)} supplies short integral source forms and bounds the archimedean norms of the untwisted lifts of their normalized images. \statementlink{thm:index-bound}{Theorem~A(b)} bounds the local lattice index of the normalized lifts. \statementlink{thm:enough-forms}{Theorem~B} shows that the images of a small number of Poincaré translates nearly span the space of good target forms. \statementlink{thm:long-forms}{Theorem~C} gives a lower bound for the covolume of a chain of theta lifts of integral forms. Applying the four estimates to the same linearly independent family, whose span has nearly full dimension in the target, yields the main inequality \eqref{eq:ch1-star} through the exact cancellation of dimension factors and the common main term.

The iteration count $m\asymp\sqrt{x}$ yields a square-root saving. Together with the prime number theorem in arithmetic progressions for a fixed modulus, which controls the error at the second scale, it turns the main inequality into a one-sided upper bound for the weighted remainder over inert prime powers. A zero-density estimate for families of quadratic twists then rules out pole cancellation. Landau's positivity criterion converts the one-sided remainder bound into nonvanishing on the right half-plane $\RePart s>1/2$. Finally, the functional equations of the Dirichlet $L$-functions and the Riemann $\zeta$-function give the critical-line conclusions and complete the proofs of the five public theorems.

\cleardoublepage
\pdfbookmark[0]{Contents}{monograph-contents}
\tableofcontents
